
\section{Minimal Tunneling Measurement Model}

We consider a minimal model of a tunneling particle coupled to a measuring device, represented by a single pointer degree of freedom. The total system is described in a two-dimensional configuration space with coordinates $(x,y)$, where $x$ denotes the particle position and $y$ the pointer position.

\subsection{Hamiltonian}

The total Hamiltonian is
\begin{equation}
\hat H =
-\frac{\hbar^2}{2m}\frac{\partial^2}{\partial x^2}
+ V(x)
-\frac{\hbar^2}{2M}\frac{\partial^2}{\partial y^2}
+ g(t)\,\hat A\,\hat P_y,
\end{equation}
where $\hat P_y = -i\hbar \frac{\partial}{\partial y}$ is the momentum conjugate to the pointer coordinate, and $g(t)$ is a time-dependent coupling.

\subsection{Barrier Potential}

We take a rectangular barrier
\begin{equation}
V(x) =
\begin{cases}
V_0, & 0 < x < a, \\
0, & \text{otherwise}.
\end{cases}
\end{equation}

\subsection{M\section{Minimal Tunneling Measurement Model}

We consider a minimal model of a tunneling particle coupled to a measuring device, represented by a single pointer degree of freedom. The total system is described in a two-dimensional configuration space with coordinates $(x,y)$, where $x$ denotes the particle position and $y$ the pointer position.

\subsection{Hamiltonian}

The total Hamiltonian is
\begin{equation}
\hat H =
-\frac{\hbar^2}{2m}\frac{\partial^2}{\partial x^2}
+ V(x)
-\frac{\hbar^2}{2M}\frac{\partial^2}{\partial y^2}
+ g(t)\,\hat A\,\hat P_y,
\end{equation}
where $\hat P_y = -i\hbar \frac{\partial}{\partial y}$ is the momentum conjugate to the pointer coordinate, and $g(t)$ is a time-dependent coupling.

\subsection{Barrier Potential}

We take a rectangular barrier
\begin{equation}
V(x) =
\begin{cases}
V_0, & 0 < x < a, \\
0, & \text{otherwise}.
\end{cases}
\end{equation}

\subsection{Measurement Coupling}

To distinguish transmitted and reflected channels, we introduce projectors $\hat \Pi_T$ and $\hat \Pi_R$ satisfying
\begin{equation}
\hat \Pi_T + \hat \Pi_R = I, \qquad \hat \Pi_T \hat \Pi_R = 0.
\end{equation}

The measured observable is defined as
\begin{equation}
\hat A = \hat \Pi_T - \hat \Pi_R,
\end{equation}
so that the interaction Hamiltonian becomes
\begin{equation}
\hat H_{\text{int}} = g(t)(\hat \Pi_T - \hat \Pi_R)\hat P_y.
\end{equation}

\subsection{Initial State}

We take a factorized initial state
\begin{equation}
\Psi(x,y,0) = \psi_{\text{in}}(x)\,\chi_0(y),
\end{equation}
where $\psi_{\text{in}}(x)$ is an incoming Gaussian wave packet
\begin{equation}
\psi_{\text{in}}(x) =
\frac{1}{(2\pi \sigma_x^2)^{1/4}}
\exp\left[-\frac{(x-x_0)^2}{4\sigma_x^2} + i k_0 x\right],
\end{equation}
and $\chi_0(y)$ is a narrow pointer wave packet
\begin{equation}
\chi_0(y) =
\frac{1}{(2\pi \sigma_y^2)^{1/4}}
\exp\left[-\frac{y^2}{4\sigma_y^2}\right].
\end{equation}

\subsection{Scattering Structure}

Under the barrier potential, the particle wave function evolves into transmitted and reflected components
\begin{equation}
\psi(x,t) = \psi_T(x,t) + \psi_R(x,t),
\end{equation}
with
\begin{equation}
\psi_T(x,t) = \int dk\, f(k)\,T(k)\,e^{ikx - iE_k t/\hbar},
\end{equation}
\begin{equation}
\psi_R(x,t) = \int dk\, f(k)\,R(k)\,e^{-ikx - iE_k t/\hbar},
\end{equation}
where $E_k = \hbar^2 k^2 / (2m)$.

\subsection{Pointer Shift}

The interaction translates the pointer depending on the scattering channel. Defining
\begin{equation}
\kappa = \int g(t)\,dt,
\end{equation}
one obtains, to leading order,
\begin{equation}
\chi_T(y) \approx \chi_0(y - \kappa), \qquad
\chi_R(y) \approx \chi_0(y + \kappa).
\end{equation}

\subsection{Branch Structure}

After the interaction, the total wave function takes the form
\begin{equation}
\Psi(x,y,t) =
\psi_T(x,t)\chi_T(y,t)
+
\psi_R(x,t)\chi_R(y,t).
\end{equation}

Defining
\begin{equation}
\Psi_T = \psi_T(x,t)\chi_T(y,t), \qquad
\Psi_R = \psi_R(x,t)\chi_R(y,t),
\end{equation}
we have
\begin{equation}
\Psi = \Psi_T + \Psi_R.
\end{equation}

When the pointer packets are well separated,
\begin{equation}
\chi_T(y)\chi_R(y) \approx 0,
\end{equation}
and the two branches are effectively non-overlapping in configuration space.

\subsection{Conditional Wave Function}

The conditional wave function of the particle is defined as
\begin{equation}
\psi_{\text{cond}}(x,t) = \Psi(x, Y(t), t),
\end{equation}
where $Y(t)$ is the actual pointer position.

If $Y(t)$ lies in the support of $\chi_T$, then
\begin{equation}
\psi_{\text{cond}}(x,t) \propto \psi_T(x,t),
\end{equation}
while if $Y(t)$ lies in the support of $\chi_R$, one obtains
\begin{equation}
\psi_{\text{cond}}(x,t) \propto \psi_R(x,t).
\end{equation}

\subsection{Guidance Equations}

The Bohmian trajectories are given by
\begin{equation}
\dot X(t) = \frac{\hbar}{m}\,\mathrm{Im}\left[\frac{\partial_x \Psi}{\Psi}\right],
\end{equation}
\begin{equation}
\dot Y(t) = \frac{\hbar}{M}\,\mathrm{Im}\left[\frac{\partial_y \Psi}{\Psi}\right],
\end{equation}
evaluated at $(x,y) = (X(t),Y(t))$.

\subsection{Interpretation}

The total wave function remains a superposition of transmitted and reflected branches, but the actual configuration $(X(t),Y(t))$ lies in one branch. The conditional wave function therefore reduces to the corresponding component, yielding an effective collapse without modifying the unitary dynamics.

\subsection{Structural Summary}

This model illustrates the key features of measurement in Bohmian mechanics:

\begin{itemize}
\item The full dynamics is governed by a two-dimensional Schrödinger equation.
\item The wave function develops a branch structure correlated with the pointer.
\item Particle motion is determined by the conditional (3D) wave function.
\item The evolution of this wave function depends on the full configuration-space structure.
\end{itemize}easurement Coupling}

To distinguish transmitted and reflected channels, we introduce projectors $\hat \Pi_T$ and $\hat \Pi_R$ satisfying
\begin{equation}
\hat \Pi_T + \hat \Pi_R = I, \qquad \hat \Pi_T \hat \Pi_R = 0.
\end{equation}

The measured observable is defined as
\begin{equation}
\hat A = \hat \Pi_T - \hat \Pi_R,
\end{equation}
so that the interaction Hamiltonian becomes
\begin{equation}
\hat H_{\text{int}} = g(t)(\hat \Pi_T - \hat \Pi_R)\hat P_y.
\end{equation}

\subsection{Initial State}

We take a factorized initial state
\begin{equation}
\Psi(x,y,0) = \psi_{\text{in}}(x)\,\chi_0(y),
\end{equation}
where $\psi_{\text{in}}(x)$ is an incoming Gaussian wave packet
\begin{equation}
\psi_{\text{in}}(x) =
\frac{1}{(2\pi \sigma_x^2)^{1/4}}
\exp\left[-\frac{(x-x_0)^2}{4\sigma_x^2} + i k_0 x\right],
\end{equation}
and $\chi_0(y)$ is a narrow pointer wave packet
\begin{equation}
\chi_0(y) =
\frac{1}{(2\pi \sigma_y^2)^{1/4}}
\exp\left[-\frac{y^2}{4\sigma_y^2}\right].
\end{equation}

\subsection{Scattering Structure}

Under the barrier potential, the particle wave function evolves into transmitted and reflected components
\begin{equation}
\psi(x,t) = \psi_T(x,t) + \psi_R(x,t),
\end{equation}
with
\begin{equation}
\psi_T(x,t) = \int dk\, f(k)\,T(k)\,e^{ikx - iE_k t/\hbar},
\end{equation}
\begin{equation}
\psi_R(x,t) = \int dk\, f(k)\,R(k)\,e^{-ikx - iE_k t/\hbar},
\end{equation}
where $E_k = \hbar^2 k^2 / (2m)$.

\subsection{Pointer Shift}

The interaction translates the pointer depending on the scattering channel. Defining
\begin{equation}
\kappa = \int g(t)\,dt,
\end{equation}
one obtains, to leading order,
\begin{equation}
\chi_T(y) \approx \chi_0(y - \kappa), \qquad
\chi_R(y) \approx \chi_0(y + \kappa).
\end{equation}

\subsection{Branch Structure}

After the interaction, the total wave function takes the form
\begin{equation}
\Psi(x,y,t) =
\psi_T(x,t)\chi_T(y,t)
+
\psi_R(x,t)\chi_R(y,t).
\end{equation}

Defining
\begin{equation}
\Psi_T = \psi_T(x,t)\chi_T(y,t), \qquad
\Psi_R = \psi_R(x,t)\chi_R(y,t),
\end{equation}
we have
\begin{equation}
\Psi = \Psi_T + \Psi_R.
\end{equation}

When the pointer packets are well separated,
\begin{equation}
\chi_T(y)\chi_R(y) \approx 0,
\end{equation}
and the two branches are effectively non-overlapping in configuration space.

\subsection{Conditional Wave Function}

The conditional wave function of the particle is defined as
\begin{equation}
\psi_{\text{cond}}(x,t) = \Psi(x, Y(t), t),
\end{equation}
where $Y(t)$ is the actual pointer position.

If $Y(t)$ lies in the support of $\chi_T$, then
\begin{equation}
\psi_{\text{cond}}(x,t) \propto \psi_T(x,t),
\end{equation}
while if $Y(t)$ lies in the support of $\chi_R$, one obtains
\begin{equation}
\psi_{\text{cond}}(x,t) \propto \psi_R(x,t).
\end{equation}

\subsection{Guidance Equations}

The Bohmian trajectories are given by
\begin{equation}
\dot X(t) = \frac{\hbar}{m}\,\mathrm{Im}\left[\frac{\partial_x \Psi}{\Psi}\right],
\end{equation}
\begin{equation}
\dot Y(t) = \frac{\hbar}{M}\,\mathrm{Im}\left[\frac{\partial_y \Psi}{\Psi}\right],
\end{equation}
evaluated at $(x,y) = (X(t),Y(t))$.

\subsection{Interpretation}

The total wave function remains a superposition of transmitted and reflected branches, but the actual configuration $(X(t),Y(t))$ lies in one branch. The conditional wave function therefore reduces to the corresponding component, yielding an effective collapse without modifying the unitary dynamics.

\subsection{Structural Summary}

This model illustrates the key features of measurement in Bohmian mechanics:

\begin{itemize}
\item The full dynamics is governed by a two-dimensional Schrödinger equation.
\item The wave function develops a branch structure correlated with the pointer.
\item Particle motion is determined by the conditional (3D) wave function.
\item The evolution of this wave function depends on the full configuration-space structure.
\end{itemize}
